% Приёмы, на которых стоит новая выгрузка. Каждый проверен сборкой обычным pdflatex.
% Страница 1: кривые и области. Страница 2: оси и сетка в четырёх положениях окна.
\documentclass[12pt,a4paper]{article}
\usepackage[T2A]{fontenc}
\usepackage[utf8]{inputenc}
\usepackage[english,russian]{babel}
\usepackage{amsmath}
\usepackage{pgfplots}
\pgfplotsset{compat=1.18}
\usetikzlibrary{arrows.meta}
\usetikzlibrary{patterns}
\usepgfplotslibrary{fillbetween}
\usepackage{geometry}
\geometry{margin=2cm}
\definecolor{cD}{HTML}{2F6FED}
\definecolor{cS}{HTML}{E0563B}
\definecolor{cG}{HTML}{2E9E44}
\definecolor{cK}{HTML}{6B6457}
\definecolor{cGrid}{HTML}{5C5346}
\begin{document}
\pagestyle{empty}
\begin{figure}[h]
\centering
\begin{tikzpicture}
\begin{axis}[name=panelA, at={(0pt,0pt)}, anchor=south west, width=340.5pt, height=226.5pt, scale only axis, clip=false, xmin=0, xmax=250, ymin=0, ymax=250, axis x line=middle, axis y line=middle, axis line style={black, line width=0.68pt, -{Stealth[length=5pt]}}, tick align=outside, major tick length=3pt, scaled ticks=false, xtick={50,100,150,200}, ytick={50,100,150,200}, tick label style={font=\fontsize{8}{9.2}\selectfont}, legend style={at={(0.8,0.95)}, anchor=north west, draw=cK, fill=white, fill opacity=0.9, text opacity=1, font=\fontsize{9}{10.3}\selectfont, cells={anchor=west}}]
\begin{scope}
\clip (axis cs:0,0) rectangle (axis cs:250,250);
% 1. Область: одна граница кривая, вторая число. Формула и два угла по прямой границе.
\addplot[draw=none, fill=cD, fill opacity=0.16, domain=0:31.9113, samples=81, forget plot] {2000 / (x + 10) - 10} -- (axis cs:31.9113,37.7287) -- (axis cs:0,37.7287) -- cycle;
\addplot[draw=none, fill=cS, fill opacity=0.16, domain=0:31.9113, samples=41, forget plot] {x^2 / 80 + 5} -- (axis cs:31.9113,17.7287) -- (axis cs:0,17.7287) -- cycle;
% 2. Область: обе границы числа. Прямоугольник.
\fill[cG, opacity=0.22] (axis cs:0,17.7287) rectangle (axis cs:31.9113,37.7287);
% 3. Область: обе границы кривые. Два невидимых пути с именами и fill between.
\addplot[draw=none, name path=areaA1hi, domain=31.9113:43.0106, samples=41, forget plot] {2000 / (x + 10) - 10};
\addplot[draw=none, name path=areaA1lo, domain=31.9113:43.0106, samples=41, forget plot] {x^2 / 80 + 5};
\addplot[cK, opacity=0.28, forget plot] fill between[of=areaA1hi and areaA1lo];
% 4. Кривые формулой: отрезок точный, отсчётов столько, сколько нужно кривизне.
\addplot[cD, line width=1.13pt, domain=0:190, samples=321, forget plot] {2000 / (x + 10) - 10};
\addplot[cS, line width=1.13pt, domain=0:140, samples=41, forget plot] {x^2 / 80 + 5};
% 5. Кусочная кривая: по куску на участок, граница точная, прямой хватает двух отсчётов.
\addplot[cG, line width=1.13pt, domain=0:40, samples=2, forget plot] {0.5 * x + 120};
\addplot[cG, line width=1.13pt, domain=40:90, samples=2, forget plot] {2 * x + 60};
% 6. Запись «по вертикали», Q = 240 - 2P: параметрически, переменная t.
\addplot[cK, line width=1.13pt, dash pattern=on 3pt off 2.4pt, domain=0:120, samples=2, variable=\t, forget plot] ({240 - 2 * t}, {t});
\addplot[cK, line width=0.9pt, domain=5:240, samples=161, variable=\t, forget plot] ({3000 / (t + 10)}, {t});
% 7. Граница, прижатая к нулю: max(0, MC). Излом закрывает сама кривая.
\addplot[draw=none, fill=cK, fill opacity=0.2, domain=150:230, samples=81, forget plot] {max(0, 1.5 * x - 270)} -- (axis cs:230,0) -- (axis cs:150,0) -- cycle;
\addplot[black, line width=0.9pt, domain=150:250, samples=161, forget plot] {max(0, 1.5 * x - 270)};
% 8. Условие с «не число» вне участка и защита от ухода далеко за окно.
\addplot[cD, line width=0.9pt, dash pattern=on 3pt off 2.4pt, domain=100:240, samples=81, restrict y to domain*=-75:325, forget plot] {((x >= 100 && x < 200) ? 220 - 0.5 * x : nan)};
% 9. Штриховка вместо заливки.
\fill[pattern=north east lines, pattern color=cG] (axis cs:200,150) rectangle (axis cs:240,200);
\end{scope}
% 10. Точки, пунктиры, подписи: точные координаты; сдвиг от точки привязки в пунктах.
\draw[cK, line width=0.45pt, dash pattern=on 2.4pt off 1.8pt] (axis cs:31.9113,37.7287) -- (axis cs:31.9113,0);
\draw[cK, line width=0.45pt, dash pattern=on 2.4pt off 1.8pt] (axis cs:31.9113,37.7287) -- (axis cs:0,37.7287);
\fill[cD] (axis cs:31.9113,37.7287) circle[radius=2.4pt];
\draw[cS, line width=0.68pt, fill=white] (axis cs:31.9113,17.7287) circle[radius=2.4pt];
\node[anchor=east, inner sep=0pt, font=\fontsize{8}{9.2}\selectfont\boldmath, xshift=-6pt] at (axis cs:0,37.7287) {$37{,}73_{b}$};
\node[anchor=north, inner sep=0pt, font=\fontsize{8}{9.2}\selectfont, yshift=-6pt] at (axis cs:31.9113,0) {$31{,}91_{1}$};
\node[anchor=north east, inner sep=0pt, xshift=-4pt, yshift=-4pt, font=\fontsize{8}{9.2}\selectfont] at (axis cs:0,0) {$0$};
% 11. Легенда своими строками: у каждого \addplot стоит forget plot.
\addlegendimage{legend image code/.code={\fill[cD, opacity=0.35] (0cm,-0.09cm) rectangle (0.42cm,0.11cm);}} \addlegendentry{$\mathrm{CS}$}
\addlegendimage{legend image code/.code={\fill[cS, opacity=0.35] (0cm,-0.09cm) rectangle (0.42cm,0.11cm);}} \addlegendentry{$\mathrm{PS}$}
\end{axis}
\end{tikzpicture}
\caption{Кривые формулами, области границами}
\end{figure}
\newpage
\begin{figure}[h]
\centering
\begin{tikzpicture}
% А. Обе оси в нуле. Линия сетки на краю окна (250) без числа и засечки: extra ticks. Мелкая сетка явным списком.
\begin{axis}[name=panelA, at={(0pt,0pt)}, anchor=south west, width=200pt, height=140pt, scale only axis, clip=false, xmin=0, xmax=250, ymin=0, ymax=250, axis x line=middle, axis y line=middle, axis line style={black, line width=0.68pt, -{Stealth[length=5pt]}}, tick align=outside, major tick length=3pt, scaled ticks=false, xtick={50,100,150,200}, ytick={50,100,150,200}, extra x ticks={250}, extra x tick labels={}, extra y ticks={250}, extra y tick labels={}, extra tick style={grid=major, major tick length=0pt}, minor xtick={10,20,30,40,60,70,80,90}, minor ytick={10,20,30,40}, minor tick length=0pt, grid=both, grid style={line width=0.4pt, cGrid, opacity=0.13}, minor grid style={line width=0.4pt, cGrid, opacity=0.06}, tick label style={font=\fontsize{8}{9.2}\selectfont}, xlabel={$Q$}, xlabel style={at={(current axis.right of origin)}, anchor=west}, ylabel={$P$}, ylabel style={at={(current axis.above origin)}, anchor=south}]
\begin{scope}
\clip (axis cs:0,0) rectangle (axis cs:250,250);
\addplot[cD, line width=1.13pt, domain=0:190, samples=321, forget plot] {2000 / (x + 10) - 10};
\end{scope}
\end{axis}
% Б. Горизонтальной оси нет (ноль по вертикали уехал за кадр), вертикальная есть. Сетка остаётся в обе стороны.
\begin{axis}[name=panelB, at={(250pt,0pt)}, anchor=south west, width=200pt, height=140pt, scale only axis, clip=false, xmin=0, xmax=100, ymin=30, ymax=90, axis x line=bottom, axis y line=middle, axis line style={black, line width=0.68pt, -{Stealth[length=5pt]}}, x axis line style={draw=none}, tick align=outside, major tick length=3pt, scaled ticks=false, xtick=\empty, ytick={40,50,60,70,80}, extra x ticks={20,40,60,80,100}, extra x tick labels={}, extra y ticks={90}, extra y tick labels={}, extra tick style={grid=major, major tick length=0pt}, grid=major, grid style={line width=0.4pt, cGrid, opacity=0.13}, tick label style={font=\fontsize{8}{9.2}\selectfont}]
\begin{scope}
\clip (axis cs:0,30) rectangle (axis cs:100,90);
\addplot[cS, line width=1.13pt, domain=0:100, samples=2, forget plot] {x};
\end{scope}
\end{axis}
% В. Осей нет вовсе (окно сдвинуто), сетка есть.
\begin{axis}[name=panelC, at={(0pt,-190pt)}, anchor=south west, width=200pt, height=140pt, scale only axis, clip=false, xmin=20, xmax=80, ymin=30, ymax=90, axis x line=bottom, axis y line=left, axis line style={draw=none}, tick align=outside, major tick length=3pt, scaled ticks=false, xtick=\empty, ytick=\empty, extra x ticks={30,40,50,60,70,80}, extra x tick labels={}, extra y ticks={40,50,60,70,80,90}, extra y tick labels={}, extra tick style={grid=major, major tick length=0pt}, grid=major, grid style={line width=0.4pt, cGrid, opacity=0.13}]
\begin{scope}
\clip (axis cs:20,30) rectangle (axis cs:80,90);
\addplot[cD, line width=1.13pt, domain=20:80, samples=2, forget plot] {100 - x};
\end{scope}
\end{axis}
% Г. Сетка выключена, оси в нуле, подписи делений свои (минус и запятая как на экране).
\begin{axis}[name=panelD, at={(250pt,-190pt)}, anchor=south west, width=200pt, height=140pt, scale only axis, clip=false, xmin=-8, xmax=8, ymin=-5, ymax=12, axis x line=middle, axis y line=middle, axis line style={black, line width=0.68pt, -{Stealth[length=5pt]}}, tick align=outside, major tick length=3pt, scaled ticks=false, xtick={-6,-4,-2,2,4,6}, xticklabels={{$-6$},{$-4$},{$-2$},{$2$},{$4$},{$6$}}, ytick={-2.5,5,10}, yticklabels={{$-2{,}5$},{$5$},{$10$}}, tick label style={font=\fontsize{8}{9.2}\selectfont}]
\begin{scope}
\clip (axis cs:-8,-5) rectangle (axis cs:8,12);
\addplot[cD, line width=1.13pt, domain=-8:8, samples=41, forget plot] {abs(x) - 2};
\end{scope}
\end{axis}
\end{tikzpicture}
\caption{Оси и сетка: что на холсте, то и на листе}
\end{figure}
\end{document}
